\label{sec:applications}

This section describes practical applications of the spatial contagion sandbox to current world challenges across four domains: public health education, K-12 STEM education, policy exploration, and software engineering research. For each domain, we present concrete use cases with target user personas demonstrating how the artifact addresses real-world needs identified in Section~\ref{sec:problem}.

\subsection{Application to Public Health Education}

Healthcare worker training in infection control and outbreak management represents a critical application domain for contagion simulation tools. Traditional training approaches rely on case studies, static diagrams, and mathematical models that may not convey the dynamic, spatial nature of disease transmission in clinical settings \cite{bennett2015simulation}.

\subsubsection{Use Case 1: Hospital Outbreak Response Training}

\textbf{Target Persona}: Dr. Sarah Chen, Infection Control Nurse at a 200-bed regional hospital, responsible for training 150 nursing staff on outbreak protocols.

\textbf{Scenario}: Dr. Chen needs to train nursing staff on identifying early warning signs of a norovirus outbreak in a long-term care ward. Using the spatial contagion sandbox, she configures a simulation reflecting her hospital's specific layout: a 30-bed ward with shared dining areas, a nurses' station, and visitor access points.

\textbf{Configuration}: The simulation parameters specify:
\begin{itemize}
    \item Agent population: 30 patients (reduced mobility), 10 nurses (high mobility), 5 visitors (intermittent presence)
    \item Transmission probability: 0.15 per contact (fecal-oral route)
    \item Environmental contamination: Surfaces retain infectivity for 4 hours
    \item Intervention: Hand hygiene compliance at 70\% baseline, improved to 90\% after training
\end{itemize}

\textbf{Training Outcome}: Nursing staff observe how a single index case leads to rapid environmental contamination, with visual feedback showing hotspots around shared facilities. The simulation demonstrates how increasing hand hygiene compliance from 70\% to 90\% reduces the reproduction number $R_t$ from 2.3 to 0.8, causing the outbreak to self-extinguish. Staff can experiment with alternative intervention strategies---isolation protocols, visitor restrictions, surface disinfection schedules---and observe outcomes in real-time.

\textbf{Value Proposition}: Unlike static training materials, the simulation provides experiential learning where staff observe cause-and-effect relationships between infection control practices and outbreak trajectories. The visual, spatial representation makes abstract epidemiological concepts (reproduction number, herd immunity threshold) tangible and memorable \cite{gaba2004future}.

\subsubsection{Use Case 2: Pandemic Preparedness Workshops}

\textbf{Target Persona}: Michael Rodriguez, Emergency Preparedness Coordinator for a state health department, conducting workshops for local health officers across 50 counties.

\textbf{Scenario}: Mr. Rodriguez conducts quarterly pandemic preparedness workshops. He needs a tool to help local health officers understand trade-offs between intervention strategies (school closures, business restrictions, mask mandates) without requiring them to learn complex modeling software.

\textbf{Application}: Using pre-configured scenario templates, workshop participants explore ``what-if'' questions:
\begin{itemize}
    \item What if we close schools but keep businesses open?
    \item What if mask compliance is 60\% vs. 80\%?
    \item What if we implement interventions 5 days vs. 10 days after the first case?
\end{itemize}

Participants compare outcomes across counties with different population densities and healthcare capacities. The zero-code configuration allows health officers to modify parameters through simple JSON files, enabling rapid scenario exploration without programming expertise.

\textbf{Value Proposition}: The simulation democratizes access to epidemiological modeling, enabling non-technical stakeholders to engage with scenario planning rather than relying solely on external consultants or centralized modeling teams \cite{halpern2020representing}.

\subsection{Application to K-12 STEM Education}

Contagion simulation provides an engaging entry point for teaching computational thinking, emergence, and complex systems to K-12 students. The visual, game-like interface makes abstract concepts accessible while aligning with Next Generation Science Standards (NGSS) for systems and system models \cite{ngss2013next}.

\subsubsection{Use Case 3: Middle School Systems Thinking Module}

\textbf{Target Persona}: Jennifer Williams, 7th-grade science teacher at a public middle school, developing a 2-week module on ``Epidemics and Emergence'' as part of the life sciences curriculum.

\textbf{Scenario}: Ms. Williams wants students to understand how individual behaviors aggregate into collective outcomes---a core concept in systems thinking. Traditional instruction uses static diagrams of SIR curves, but students struggle to connect individual actions (hand washing, staying home when sick) to population-level outcomes (flattening the curve).

\textbf{Pedagogical Design}: Students work in pairs with the contagion sandbox configured for a simplified ``classroom epidemic'' scenario:
\begin{itemize}
    \item Agent population: 25 students, 1 teacher
    \item Simplified transmission: direct contact only (no environmental persistence)
    \item Adjustable parameters: student mask-wearing (yes/no), quarantine compliance (0-100\%)
\end{itemize}

Students conduct controlled experiments varying one parameter at a time and observing outcomes. Guided inquiry questions include:
\begin{enumerate}
    \item What happens if 0\% of students wear masks? 50\%? 100\%?
    \item How does the epidemic change if sick students stay home for 5 days vs. 0 days?
    \item Can you find a combination of interventions that prevents an epidemic without requiring 100\% compliance with any single measure?
\end{enumerate}

\textbf{Learning Outcomes}: Students discover emergent phenomena: the epidemic threshold, the nonlinear relationship between compliance and outcomes, and the concept of herd immunity. The visual simulation makes these concepts concrete rather than abstract mathematical relationships.

\textbf{Value Proposition}: The contagion sandbox aligns with constructionist learning principles---students learn by constructing and manipulating models rather than passively receiving instruction \cite{papert1980mindstorms}. The zero-code interface removes barriers that would otherwise require teaching programming before exploring epidemiological concepts.

\subsubsection{Use Case 4: High School AP Biology Extension}

\textbf{Target Persona}: Dr. James Park, AP Biology teacher at a magnet high school, seeking to extend the epidemiology unit with computational modeling experiences.

\textbf{Scenario}: AP Biology standards include understanding population dynamics and mathematical models in biology \cite{collegeboard2019ap}. Dr. Park wants students to experience agent-based modeling as a complement to differential equation-based SIR models.

\textbf{Application}: Students compare three modeling approaches:
\begin{enumerate}
    \item Analytical SIR model (differential equations solved with Python)
    \item Agent-based simulation (contagion sandbox)
    \item Real-world epidemic data (CDC case reports)
\end{enumerate}

Students analyze how agent-level heterogeneity (different contact rates, spatial clustering) produces epidemic curves that deviate from the homogeneous mixing assumption in standard SIR models. This comparison deepens understanding of model assumptions and their implications.

\textbf{Value Proposition}: The contagion sandbox enables authentic computational modeling experiences without requiring advanced programming skills, making agent-based modeling accessible to high school students and supporting the transition to research-based learning \cite{wilensky2015agent}.

\subsection{Application to Policy Exploration}

Policy analysts and decision-makers require rapid prototyping environments to explore intervention scenarios before committing to resource-intensive formal modeling or policy implementation. The contagion sandbox serves as a ``sketching tool'' for policy exploration.

\subsubsection{Use Case 5: School District Reopening Planning}

\textbf{Target Persona}: Linda Thompson, School Board Member in a suburban district with 15,000 students, serving on the reopening task force during an influenza season.

\textbf{Scenario}: The school board must decide between three reopening strategies: full in-person, hybrid (50\% capacity), and fully remote. Board members have access to expert epidemiological models but lack intuition for how these translate to their specific context.

\textbf{Application}: Ms. Thompson uses the contagion sandbox to explore scenarios reflecting her district's characteristics:
\begin{itemize}
    \item Full in-person: 30 students per classroom, 180 school days
    \item Hybrid: 15 students per classroom, alternating days
    \item Remote: No in-school transmission (baseline community transmission only)
\end{itemize}

She observes how different classroom configurations affect outbreak probability and peak case loads. The visual simulation helps her communicate findings to parents and community members who lack epidemiological training.

\textbf{Value Proposition}: The contagion sandbox enables non-expert stakeholders to develop intuition for epidemiological trade-offs, supporting informed decision-making and improving communication between modelers and decision-makers \cite{saltelli2020five}.

\subsubsection{Use Case 6: Workplace Intervention Strategy}

\textbf{Target Persona}: Robert Kim, Chief Operating Officer of a 500-employee technology company, developing return-to-office policies during respiratory illness season.

\textbf{Scenario}: Mr. Kim must balance employee safety, productivity, and operational continuity. He considers interventions: improved ventilation, staggered shifts, testing requirements, and mask policies.

\textbf{Application}: Using the contagion sandbox configured for an open-plan office environment, Mr. Kim explores combinations of interventions:
\begin{itemize}
    \item Baseline: No interventions, $R_t = 2.5$
    \item Ventilation upgrade: $R_t = 1.8$
    \item Ventilation + masks: $R_t = 1.1$
    \item Ventilation + masks + staggered shifts: $R_t = 0.7$
\end{itemize}

The simulation demonstrates that achieving $R_t < 1$ requires multiple interventions---single measures are insufficient. This insight informs a layered intervention strategy communicated to employees.

\textbf{Value Proposition}: The contagion sandbox enables business leaders to explore intervention trade-offs quantitatively, supporting evidence-based policy decisions without requiring in-house epidemiological modeling expertise \cite{baker2021business}.

\subsection{Application to Software Engineering Research}

Beyond domain applications, the contagion sandbox contributes to software engineering research by demonstrating composable simulation patterns that can be reused and extended by other researchers building multi-agent systems.

\subsubsection{Composable Simulation Pattern}

The artifact demonstrates a reusable architectural pattern: encapsulating simulation logic within Vue 3 composables (\texttt{useContagionEngine.js}) that can be integrated into diverse applications without modification. This pattern addresses a gap in multi-agent system research, where simulation components are often tightly coupled to specific platforms, impeding reuse \cite{kravari2015survey}.

\textbf{Pattern Characteristics}:
\begin{itemize}
    \item \textbf{Encapsulation}: Simulation state and logic isolated from presentation layer
    \item \textbf{Configuration-Driven}: Parameters externalized to JSON/YAML, enabling domain experts to customize behavior
    \item \textbf{Event-Driven}: Simulation emits events (state changes, transmission events) that UI components subscribe to
    \item \textbf{Platform-Agnostic}: Composable can be used with any rendering framework (PixiJS, Three.js, D3.js)
\end{itemize}

\textbf{Research Contribution}: This pattern can be adapted for other agent-based simulations: traffic modeling, pedestrian dynamics, ecological systems, and social network diffusion. Researchers can reuse the architecture while substituting domain-specific rules.

\subsubsection{Use Case 7: Research Tool Extension}

\textbf{Target Persona}: Dr. Maria Santos, Computer Science PhD student researching opinion dynamics in social networks, seeking a visual simulation platform for her dissertation work.

\textbf{Scenario}: Dr. Santos wants to adapt the contagion sandbox to model opinion spread rather than disease transmission. She needs to replace the infection logic with opinion update rules while preserving the visual rendering and configuration system.

\textbf{Application}: By extending the composable architecture, Dr. Santos creates \texttt{useOpinionEngine.js} that:
\begin{itemize}
    \item Replaces SIR states with opinion states (support, oppose, neutral)
    \item Implements opinion update rules based on homophily and social influence
    \item Reuses the SpatialCanvas rendering and parameter configuration infrastructure
\end{itemize}

This extension requires modifying approximately 200 lines of domain logic while reusing 800+ lines of infrastructure code (rendering, state management, configuration parsing).

\textbf{Value Proposition}: The composable architecture reduces development time for researchers building related simulations, enabling rapid prototyping and iteration \cite{gamma1994design}. This addresses a key barrier in computational social science research: the effort required to build visual, interactive simulation platforms from scratch.

\subsection{Broader Impact and Societal Implications}

The applications described above share a common theme: democratizing access to epidemiological simulation and computational modeling. By lowering barriers to entry, the contagion sandbox has potential for broader societal impact:

\textbf{Pandemic Preparedness}: Widespread familiarity with contagion simulation concepts could improve public understanding of and compliance with public health interventions during future pandemics \cite{van2021fighting}.

\textbf{STEM Workforce Development}: Engaging K-12 students with computational modeling builds skills essential for data science, epidemiology, and computational social science careers---fields experiencing workforce shortages \cite{national2020building}.

\textbf{Evidence-Based Policymaking}: Tools that enable non-experts to explore intervention scenarios support democratic participation in policy decisions with significant societal impacts \cite{saltelli2020five}.

\textbf{Open Science}: The artifact's composable architecture and configuration-driven design align with open science principles, enabling replication, extension, and adaptation by the research community \cite{munafo2017manifesto}.

\subsection{Limitations and Appropriate Use}

While the contagion sandbox serves educational and exploratory purposes effectively, we emphasize limitations to prevent misuse:

\begin{itemize}
    \item \textbf{Not for Clinical Decision-Making}: The simulation is not validated against real-world epidemiological data and should not inform medical or public health decisions requiring certified model accuracy.
    
    \item \textbf{Simplified Dynamics}: The cellular automata approach omits realistic biological mechanisms (viral load variation, asymptomatic transmission, immunity waning) captured by more sophisticated models.
    
    \item \textbf{Scale Limitations}: Performance constraints limit simulations to small populations (10-50 agents), precluding city-scale or regional-scale modeling.
    
    \item \textbf{Pedagogical Purpose}: The primary value lies in conceptual understanding and intuition-building rather than quantitative prediction.
\end{itemize}

These limitations position the artifact as a complement to, rather than replacement for, established epidemiology platforms designed for scientific research and policy analysis.

\subsection{Summary}

The spatial contagion sandbox addresses real-world needs across four domains: (1) public health education, where it enables experiential training for healthcare workers; (2) K-12 STEM education, where it makes emergence and systems thinking accessible to students; (3) policy exploration, where it democratizes access to scenario modeling for non-expert stakeholders; and (4) software engineering research, where it demonstrates reusable composable simulation patterns. Through concrete use cases with target personas, we have demonstrated how the artifact's zero-code configuration, visual interface, and integration with multi-agent platforms translate into practical impact across educational, policy, and research contexts.
